\chapter{Worked Examples and Common Pitfalls}

\begin{remarkbox}
This appendix is designed as a bridge between reading the notes and doing field
theory independently. It collects worked calculations and common mistakes that
recur throughout classical field theory.
\end{remarkbox}

\section{Worked Example: Scalar Field Equation from the Action}

Consider the real scalar action
\[
    S[\phi]=\int d^{d+1}x\left(
        \frac{1}{2}\partial_\mu\phi\partial^\mu\phi - V(\phi)
    \right).
\]
The variation is
\[
    \delta S=\int d^{d+1}x\left(
        \partial_\mu\phi\partial^\mu\delta\phi - V'(\phi)\delta\phi
    \right).
\]
Integrating the first term by parts gives
\[
    \int d^{d+1}x\,\partial_\mu\phi\partial^\mu\delta\phi
    =-\int d^{d+1}x\,\partial_\mu\partial^\mu\phi\,\delta\phi
    +\text{boundary term}.
\]
If the boundary term vanishes, then
\[
    \delta S=-\int d^{d+1}x\,\left(\Box\phi+V'(\phi)\right)\delta\phi.
\]
Since \(\delta\phi\) is arbitrary in the bulk, the Euler--Lagrange equation is
\[
    \Box\phi+V'(\phi)=0.
\]
For
\[
    V(\phi)=\frac{1}{2}m^2\phi^2,
\]
this becomes
\[
    (\Box+m^2)\phi=0.
\]

\begin{remarkbox}
Common mistake: forgetting the boundary term. The bulk equation follows only
after the boundary term has been made harmless by boundary conditions, falloff
conditions, or an additional boundary term in the action.
\end{remarkbox}

\section{Worked Example: Noether Current of a Complex Scalar Field}

The free complex scalar field has Lagrangian density
\[
    \mathcal{L}=\partial_\mu\phi^*\partial^\mu\phi-m^2\phi^*\phi.
\]
It is invariant under the global phase transformation
\[
    \phi\longrightarrow e^{i\alpha}\phi,
    \qquad
    \phi^*\longrightarrow e^{-i\alpha}\phi^*.
\]
For infinitesimal \(\alpha\), remove the common parameter and write
\[
    \Delta\phi=i\phi,
    \qquad
    \Delta\phi^*=-i\phi^*.
\]
The derivatives of the Lagrangian are
\[
    \frac{\partial\mathcal{L}}{\partial(\partial_\mu\phi)}=\partial^\mu\phi^*,
    \qquad
    \frac{\partial\mathcal{L}}{\partial(\partial_\mu\phi^*)}=\partial^\mu\phi.
\]
Therefore the Noether current is
\[
\begin{aligned}
    J^\mu
    &=\partial^\mu\phi^*(i\phi)+\partial^\mu\phi(-i\phi^*) \\
    &=i(\phi\partial^\mu\phi^*-\phi^*\partial^\mu\phi).
\end{aligned}
\]
On shell,
\[
    \partial_\mu J^\mu=0.
\]
The corresponding charge is
\[
    Q=\int d^d x\,J^0.
\]

\begin{remarkbox}
Common mistake: confusing a global phase symmetry with a gauge symmetry. A global
symmetry has a constant parameter and usually gives a conserved Noether charge. A
gauge symmetry has an arbitrary function as parameter and represents redundancy
of description; its charges require more careful boundary analysis.
\end{remarkbox}

\section{Worked Example: Maxwell Equation from the Action}

Start with
\[
    S[A]=\int d^4x\left(-\frac{1}{4}F_{\mu\nu}F^{\mu\nu}-J_\mu A^\mu\right),
\]
where
\[
    F_{\mu\nu}=\partial_\mu A_\nu-\partial_\nu A_\mu.
\]
The variation of the field strength is
\[
    \delta F_{\mu\nu}=\partial_\mu\delta A_\nu-\partial_\nu\delta A_\mu.
\]
Using antisymmetry of \(F^{\mu\nu}\),
\[
    \delta\left(-\frac{1}{4}F_{\mu\nu}F^{\mu\nu}\right)
    =-F^{\mu\nu}\partial_\mu\delta A_\nu.
\]
Integrating by parts gives
\[
    \delta S=\int d^4x\left(\partial_\mu F^{\mu\nu}-J^\nu\right)\delta A_\nu
    +\text{boundary term}.
\]
Thus the field equation is
\[
    \partial_\mu F^{\mu\nu}=J^\nu.
\]

Gauge invariance of the source term requires
\[
    \partial_\mu J^\mu=0.
\]
Indeed, under \(A_\mu\to A_\mu+\partial_\mu\chi\),
\[
    \delta S_{\mathrm{int}}
    =-\int d^4x\,J^\mu\partial_\mu\chi
    =\int d^4x\,\chi\partial_\mu J^\mu+\text{boundary term}.
\]

\begin{remarkbox}
Common mistake: treating the potential \(A_\mu\) as directly observable. The
field strength \(F_{\mu\nu}\) is gauge invariant in Abelian theory; the potential
contains redundancy.
\end{remarkbox}

\section{Worked Example: Why the Dirac Equation Squares to Klein--Gordon}

Let the free Dirac equation be
\[
    (i\gamma^\mu\partial_\mu-m)\psi=0.
\]
Apply the operator \((i\gamma^\nu\partial_\nu+m)\) from the left:
\[
    (i\gamma^\nu\partial_\nu+m)(i\gamma^\mu\partial_\mu-m)\psi=0.
\]
Expanding gives
\[
    -\gamma^\nu\gamma^\mu\partial_\nu\partial_\mu\psi-m^2\psi=0,
\]
because the cross terms cancel. Since \(\partial_\nu\partial_\mu\) is symmetric
in \(\mu,\nu\), only the anticommutator part of \(\gamma^\nu\gamma^\mu\) matters:
\[
    \gamma^\nu\gamma^\mu\partial_\nu\partial_\mu
    =\frac{1}{2}\{\gamma^\nu,\gamma^\mu\}\partial_\nu\partial_\mu
    =\eta^{\nu\mu}\partial_\nu\partial_\mu
    =\Box.
\]
Therefore
\[
    (\Box+m^2)\psi=0.
\]
Each component of a free Dirac spinor satisfies the Klein--Gordon equation, but
the Dirac equation contains additional first-order spinor information.

\begin{remarkbox}
Common mistake: saying that the Dirac equation is simply the square root of the
Klein--Gordon equation and stopping there. It is more than that: it also carries
a spinor representation of the Lorentz group and a conserved positive Dirac
current.
\end{remarkbox}

\section{Common Pitfall: Equation Without Model}

An equation becomes a physical model only after its variables and assumptions are
specified. The expression
\[
    \partial_t^2\phi-c^2\nabla^2\phi=0
\]
may describe a string, sound, a scalar perturbation, or an idealized relativistic
wave. The formal equation alone does not tell us what \(\phi\) means, what
\(c\) represents, or what boundary conditions are allowed.

A safe reading checklist is:
\begin{enumerate}
    \item identify the field;
    \item identify the domain;
    \item identify the parameters and dimensions;
    \item identify initial and boundary data;
    \item identify the symmetries;
    \item identify the conserved quantities;
    \item identify what is gauge and what is physical.
\end{enumerate}

\section{Common Pitfall: Sign Conventions}

Metric signature changes signs. Fourier conventions change signs. Gauge-coupling
conventions change signs. Gamma-matrix conventions change signs. The same theory
can look different in two books because the conventions differ.

For these notes the default convention is
\[
    \eta_{\mu\nu}=\operatorname{diag}(+1,-1,-1,-1),
    \qquad
    \Box=\partial_t^2-\nabla^2.
\]
With this convention the Klein--Gordon equation is
\[
    (\Box+m^2)\phi=0.
\]
Before comparing formulas with another source, check the signature and Fourier
convention first.

\section{Common Pitfall: Local Conservation vs Conserved Charge}

A local conservation law has the form
\[
    \partial_\mu J^\mu=0.
\]
In space-plus-time notation this is
\[
    \partial_t\rho+\nabla\cdot\vect{J}=0.
\]
The charge in a region \(R\) is
\[
    Q_R(t)=\int_R d^d x\,\rho(t,\vect{x}).
\]
Its time derivative is
\[
    \frac{dQ_R}{dt}=-\int_{\partial R}\vect{J}\cdot d\vect{S}.
\]
Thus a locally conserved current does not imply that the charge inside every
finite region is constant. It implies that charge changes only by flux through
the boundary. Total charge is conserved only when the boundary flux vanishes or
when the region is all space with suitable falloff.

\section{Common Pitfall: Gauge Choice vs Physics}

A gauge condition is a choice of description, not an additional physical law. For
example,
\[
    \partial_\mu A^\mu=0
\]
is the Lorenz gauge, while
\[
    \nabla\cdot\vect{A}=0
\]
is the Coulomb gauge. These conditions restrict the representative \(A_\mu\), not
the physical electromagnetic field itself.

Gauge-invariant or gauge-covariant quantities carry the physical content. In
Abelian electromagnetism, \(F_{\mu\nu}\) is gauge invariant. In non-Abelian gauge
theory, \(F_{\mu\nu}\) transforms covariantly, and gauge-invariant quantities are
constructed by traces, invariant inner products, or suitable nonlocal objects.

\section{Common Pitfall: Classical Spinor Field vs Fermion}

The Dirac equation can be treated as a classical relativistic field equation, but
fermions are quantum objects. The physical spin-statistics connection and
anticommutation relations belong to quantum field theory. In a classical action
formalism, spinor components are often treated as Grassmann-valued variables to
prepare for quantization.

The safe statement is this: the classical Dirac action, Lorentz covariance,
current conservation, and gauge coupling can be studied classically. The full
interpretation of the Dirac field as a fermion field is quantum.

\begin{summarybox}
Most technical errors in classical field theory come from a small number of
sources: ignoring boundary terms, confusing global and gauge symmetries, changing
sign conventions without noticing, mistaking gauge choices for physics, and
forgetting that an equation is not a complete model until its variables,
assumptions, and boundary data have been specified.
\end{summarybox}
